\begin{center}
\hypertarget{native-033-article015}{}\pdfbookmark[2]{15.}{article-015}15.
\end{center}
Quum anomalia media $M$, ut vidimus, per $v$ et $\varphi$ complete determinata sit, sicuti $v$ per $M$ et $\varphi$, patet, si omnes tres quantitates simul ut variabiles spectentur, inter ipsarum variationes differentiales aequationem conditionalem locum habere debere, cuius investigatio haud superflua erit. Differentiando \SourceLeaf{034}primo aequationem VII art. 8, prodit $\frac{\mathrm{d}E}{\sin E}=\frac{\mathrm{d}v}{\sin v}-\frac{\mathrm{d}\varphi}{\cos\varphi}$; differentiando perinde aequationem XII, fit $\mathrm{d}M=(1-e\cos E)\mathrm{d}E-\sin E\cos\varphi\,\mathrm{d}\varphi$. Eliminando ex his aequationibus differentialibus $\mathrm{d}E$, obtinemus
\[
\mathrm{d}M=\frac{\sin E(1-e\cos E)}{\sin v}\,\mathrm{d}v-\left(\sin E\cos\varphi+\frac{\sin E(1-e\cos E)}{\cos\varphi}\right)\mathrm{d}\varphi
\]
sive substituendo pro $\sin E$, $1-e\cos E$ valores suos ex aequatt. VIII, III
\[
\mathrm{d}M=\frac{r r}{a a\cos\varphi}\,\mathrm{d}v-\frac{r(r+p)\sin v}{a a\cos\varphi^2}\,\mathrm{d}\varphi
\]
sive denique, exprimendo utrumque coëfficientem per $v$ et $\varphi$ tantum,
\[
\mathrm{d}M=\frac{\cos\varphi^3}{(1+e\cos v)^2}\,\mathrm{d}v-\frac{(2+e\cos v)\sin v\cos\varphi^2}{(1+e\cos v)^2}\,\mathrm{d}\varphi.
\]
Vice versa considerando $v$ tamquam functionem quantitatum $M$, $\varphi$, aequatio hancce formam obtinet:
\[
\mathrm{d}v=\frac{a a\cos\varphi}{r r}\,\mathrm{d}M+\frac{(2+e\cos v)\sin v}{\cos\varphi}\,\mathrm{d}\varphi
\]
sive introducendo $E$ pro $v$\,\raisebox{-.35ex}{.}
\[
\mathrm{d}v=\frac{a a\cos\varphi}{r r}\,\mathrm{d}M+\frac{a a}{r r}(2-e\cos E-e e)\sin E\,\mathrm{d}\varphi.
\]

\begin{center}
\hypertarget{native-034-article016}{}\pdfbookmark[2]{16.}{article-016}16.
\end{center}
Radius vector $r$ per $v$ et $\varphi$ vel per $M$ et $\varphi$ plene nondum determinatus est, sed insuper a $p$ vel $a$ pendet; constabit igitur eius differentiale tribus membris. Per differentiationem aequationis II art. 8 nanciscimur
\[
\frac{\mathrm{d}r}{r}=\frac{\mathrm{d}p}{p}+\frac{e\sin v}{1+e\cos v}\,\mathrm{d}v-\frac{\cos\varphi\cos v}{1+e\cos v}\,\mathrm{d}\varphi.
\]
Statuendo hic $\frac{\mathrm{d}p}{p}=\frac{\mathrm{d}a}{a}-2\operatorname{tang}\varphi\,\mathrm{d}\varphi$ (quod sequitur e differentiatione aequ. I), exprimendoque secundum art. praec. $\mathrm{d}v$ per $\mathrm{d}M$ et $\mathrm{d}\varphi$, prodit post debitas reductiones
\[
\frac{\mathrm{d}r}{r}=\frac{\mathrm{d}a}{a}+\frac{a}{r}\operatorname{tang}\varphi\sin v\,\mathrm{d}M-\frac{a}{r}\cos\varphi\cos v\,\mathrm{d}\varphi
\]
sive
\[
\mathrm{d}r=\frac{r}{a}\,\mathrm{d}a+a\operatorname{tang}\varphi\sin v\,\mathrm{d}M-a\cos\varphi\cos v\,\mathrm{d}\varphi.
\]

Ceterum hae formulae, sicut eae quas in art. praec. evolvimus, suppositioni innituntur, $v$, $\varphi$ et $M$ sive potius $\mathrm{d}v$, $\mathrm{d}\varphi$ et $\mathrm{d}M$ in partibus radii \SourceLeaf{035}exprimi. Quodsi igitur variationes angulorum $v$, $\varphi$, $M$ in secundis exprimere placet: vel eas formularum partes quae $\mathrm{d}v$, $\mathrm{d}\varphi$ aut $\mathrm{d}M$ implicant, per $206264{,}8$ dividere oportet, vel eas, quae continent $\mathrm{d}r$, $\mathrm{d}p$ aut $\mathrm{d}a$, per eundem numerum multiplicare. Formulae igitur art. praec., quae hoc respectu sunt homogeneae, mutatione opus non habebunt.
